%! Piecewise Linear Learning Functions — Unit 8, Lesson 8.1 — Exit Ticket
\documentclass[10pt]{article}
\usepackage{linalg-article}
\usepackage{linalg-boxes}
\newcommand{\TallMath}[1]{$\displaystyle #1\rule[-1.4em]{0pt}{3.2em}$}
\newcommand{\relu}{\mathrm{ReLU}}

\begin{document}

\pageheader{Unit 8, Lesson 8.1}{Exit Ticket}
\namedateperiod

\vspace{0.15in}

No notes. Show your reasoning.

\begin{enumerate}[leftmargin=1.8em, itemsep=16pt]
  \item \textbf{Build the shape.} Let $F(x)=\relu(x)-2\,\relu(x-2)+\relu(x-4)$. Compute
        $F(0)$, $F(2)$, $F(4)$, and $F(6)$, and describe the shape the graph makes.
        \vspace{1.5cm}
  \item \textbf{Folds and pieces.} A piecewise-linear function is built from $4$ ReLUs. How many folds does
        it have, and how many straight pieces?
        \vspace{1.1cm}
  \item \textbf{Explain.} In one sentence, why can a sum of ReLUs bend to fit rising-and-falling data when a
        single straight line cannot?
        \writelines{2}
\end{enumerate}

\end{document}
